\SourceLeaf{21}

The Observatio can easily be derived from the Criterien. For\,example, for 3 to be a cubic residue,
the condition $N\equiv 0\;(\mathrm{mod.}\ 3)$ was just obtained and is likewise stated under [2]. If $M$ and $N$
are even, put
\[ M=2x,\qquad N=6y, \]
so that $4n=M^2+27N^2$ implies that $n$ is representable by the quadratic form of determinant $-243$
\[ n=x^2+243y^2 \]
with integers $x, y$. Conversely, the possibility of such a representation of $n$
implies the validity of the congruence $N\equiv 0\;(\mathrm{mod.}\ 3)$. If $M$ and $N$ are odd, choose the appropriate sign and put
\[ M\pm\dfrac{N}{3}=4x,\qquad \mp\dfrac{N}{3}=y, \]
which gives the representation
\[ n=4x^2+2xy+61y^2. \]
Conversely, this formula would imply $N\equiv 0\;(\mathrm{mod.}\ 3)$, so that 3 is a cubic residue
modulo $n$ if and only if $n$ is representable by one of the forms $(1,0,243)$, $(4,1,61)$.

Mr. R. \textsc{Dedekind}, to whom I owe several very helpful conversations on the subjects in question,
drew my attention to the fact that, in the case where
3 is a cubic nonresidue modulo $n$, \textsc{Gauss}'s statements are not quite complete. There are altogether nine classes of primitive forms of the first kind of determinant $D=-243$, whose reduced forms are the following
\[ (1,0,243),\quad (4,\pm1,61),\quad (13,\pm2,19),\quad (7,\pm3,36),\quad (9,\pm3,28). \]
In [1] the forms $(13,\pm2,19)$ are missing. The theorem must therefore state that 3 is a cubic
nonresidue modulo $n$ if and only if $n$ admits a representation in one of the forms
\[ n=13xx+4xy+19yy, \]
\[ n=7xx+6xy+36yy, \]
\[ n=9xx+6xy+28yy. \]
These representations can always be obtained by choosing the appropriate sign in the following transformations
\[ \text{I. If } M\equiv N\equiv 1\;(\mathrm{mod.}\ 2),\quad M\pm 7N\equiv 0\;(\mathrm{mod.}\ 12), \]
\[ M=5x+7y,\qquad \pm N=x-y. \]
\[ \text{II. If } M\equiv N\equiv 1\;(\mathrm{mod.}\ 2),\quad M\pm 7N\equiv 6\;(\mathrm{mod.}\ 12), \]
\[ M=x+6y,\qquad \mp N=x. \]
\[ \text{III. If } M\equiv N\equiv 0\;(\mathrm{mod.}\ 2),\quad M\pm N\equiv 0\;(\mathrm{mod.}\ 3), \]
\[ M=6x+2y,\qquad \mp N=2y. \]
\par\noindent\hspace*{\parindent}The statements about the number 5 given under 1 follow in a similar way from the residue criterion
$MN\equiv 0\;(\mathrm{mod.}\ 5)$.

It should also be noted that the original of note [1] contains a table extending to 223 of the primes $p=3n+1$ modulo which the numbers 2, 3, 5, 7 are cubic residues.

\hfill \textsc{Fricke.}

\endgroup
